\documentclass[12pt,a4paper]{article}
\usepackage[margin=2.4cm]{geometry}
\usepackage{amsmath,amssymb,bm,mathtools}
\usepackage{fontspec}
\setmainfont{TeX Gyre Termes}
\setsansfont{TeX Gyre Heros}
\setmonofont{Latin Modern Mono}[Scale=0.9]
\usepackage{graphicx,float,caption,subcaption,booktabs,longtable,array,multirow,tabularx}
\usepackage{xcolor}
\usepackage{setspace}
\usepackage{lineno}
\usepackage{microtype}
\usepackage{xurl}
\usepackage{siunitx}
\usepackage{enumitem}
\usepackage{fancyhdr}
\usepackage{titlesec}
\titleformat{\paragraph}[runin]{\normalfont\itshape}{}{0pt}{}
\titlespacing*{\paragraph}{0pt}{1.1ex plus .2ex}{0.7em}
\usepackage[authoryear,round,semicolon,sort&compress]{natbib}
\definecolor{lk}{RGB}{0,70,140}
\usepackage[colorlinks=true,linkcolor=lk,citecolor=lk,urlcolor=lk,breaklinks=true,pdfborder={0 0 0}]{hyperref}
\usepackage[nameinlink,capitalise]{cleveref}
\sisetup{detect-all,separate-uncertainty=false}
\setstretch{1.5}
\setlength{\parskip}{3pt}
\captionsetup{font=small,labelfont=bf,justification=centering,singlelinecheck=true}
\titleformat{\section}{\large\bfseries}{\thesection.}{0.5em}{}
\titleformat{\subsection}{\normalsize\bfseries}{\thesubsection}{0.5em}{}
\titleformat{\subsubsection}{\normalsize\itshape}{\thesubsubsection}{0.5em}{}
\pagestyle{fancy}\fancyhf{}\fancyfoot[C]{\small\thepage}\renewcommand{\headrulewidth}{0pt}
\setlength{\headheight}{14pt}
\bibliographystyle{plainnathl}
\newcommand{\doihl}[1]{doi: \href{https://doi.org/#1}{\nolinkurl{#1}}}
\newcommand{\dd}{\mathrm{d}}
\newcommand{\ET}{\mathit{ET}}
\newcommand{\Ent}{\mathit{ET}_{\mathrm{nt}}}
\newcommand{\Ks}{K_{\mathrm{s}}}
\newcommand{\Se}{S_{\mathrm{e}}}
\newcommand{\TAW}{\mathit{TAW}}
\newcommand{\repourl}{https://github.com/halilunaziru73-creator/Coupled-Water-Fluxes-Pore-Scale-Dynamics-Precision-Irrigation-Mulching-Feedback-Loops}
\newcommand{\repo}{\href{\repourl}{repository}}
\newcommand{\datalink}{\href{\repourl}{Data}}
\newcommand{\codelink}{\href{\repourl/tree/main/Manuscripts}{Code}}
\newcommand{\softlink}{\href{\repourl/tree/main/Software/SWAFL_Desktop}{Software}}
\newcommand{\manuallink}{\href{\manualurl}{Manual}}
\newcommand{\downloadlink}{\href{\exeurl}{SWAFL.exe}}
\newcommand{\figdir}{../figures}
\newcommand{\manualurl}{https://github.com/halilunaziru73-creator/Coupled-Water-Fluxes-Pore-Scale-Dynamics-Precision-Irrigation-Mulching-Feedback-Loops/blob/main/Software/SWAFL_Desktop/docs/SWAFL_Manual.pdf}
\newcommand{\exeurl}{https://github.com/halilunaziru73-creator/Coupled-Water-Fluxes-Pore-Scale-Dynamics-Precision-Irrigation-Mulching-Feedback-Loops/releases/latest/download/SWAFL.exe}
\newcommand{\authorblock}{%
\begin{center}
{\large\textbf{Naziru Halilu}}$^{1,2,*}$\\[3pt]
{\small $^{1}$ Agricultural University of Athens, Athens, Greece\\
$^{2}$ Public University of Navarre, Pamplona, Spain\\
ORCID: \href{https://orcid.org/0009-0006-2324-9599}{0009-0006-2324-9599}\\
$^{*}$ Correspondence: \href{mailto:halilunaziru73@gmail.com}{halilunaziru73@gmail.com}}
\end{center}}
\newcommand{\codeavail}{%
\datalink{} (the complete Daisy ensemble outputs), \codelink{} (analysis and figure scripts) and \softlink{} (the SWAFL Desktop source) are in the project repository; the \manuallink{} (93 pages) can be reviewed or downloaded as a PDF. The dataset (all 10\,190 runs, raw and derived) is archived on Zenodo, \doihl{10.5281/zenodo.23043626}, and the ensemble table used here is also stored in the repository as \texttt{Manuscripts/reproducibility/data/daisy\_ensemble.csv.gz}. The folder \texttt{Manuscripts/reproducibility} is a self-contained package: pinned software versions (Python 3.13.16, NumPy 2.5.3, SciPy 1.18.1, pandas 3.0.5, scikit-learn 1.9.1, Matplotlib 3.11.2, Node.js 22.22.0), fixed random seeds in every stochastic step, SHA-256 checksums, and one command (\texttt{run\_all.sh}) that regenerates all results and figures from the ensemble and then compares every regenerated value with the archived result and with the headline numbers printed in the manuscripts. Each table and figure names the script that produced it (Appendix~\ref{app:repro}). Daisy itself was not re-run for this package, no field measurements are used, and three small intermediate arrays and the interface screenshots are archived rather than regenerated; these limits are listed in \texttt{README\_REPRODUCE.md}.}

\usepackage{pdflscape}
\setlength{\emergencystretch}{3em}

\hypersetup{pdftitle={Auditing Water-Balance Closure in a Soil-Plant Model with Interpretable Machine Learning},pdfauthor={Naziru Halilu}}
\begin{document}
\linenumbers
\begin{titlepage}
\vspace*{-1.2cm}
\begin{center}{\Large\bfseries Auditing Water-Balance Closure in a Soil--Plant Model with Interpretable Machine Learning\par}\end{center}
\vspace{0.4cm}
\authorblock
\vspace{0.4cm}

{\small\setlength{\parindent}{0pt}\setlength{\parskip}{0pt}
\begin{center}\textbf{\normalsize Abstract}\end{center}\vspace{-0.4em}
\noindent Scientific machine learning promises equations discovered from data, yet few discovery pipelines are first shown to recover a known answer, and fewer ask whether the quantities they are given close in the first place. A discovered relation is only as credible as the controls it has passed. A discovery protocol is developed in which a blind linear-invariant search must return the water balance before anything else is believed, with a permutation null and resampling to calibrate its uncertainty, and in which every partition used as a target is audited for closure. The protocol is applied to a large designed ensemble of simulations from a widely used soil--plant model, across experiments and soils that are withheld in turn. The search recovers the water balance with quantified uncertainty, and the closure audit then shows what the reported components leave out: an unattributed evaporative flux that is a substantial share of evapotranspiration, and two experiments in which stored water does not close. Because this flux changes the partition fractions that are used to compare management treatments, its identification matters to users of the model and of emulators built on it. A compact law with three terms predicts evapotranspiration in withheld experiments and withheld soils as accurately as gradient boosting, whereas drainage resists both, which locates the information boundary of annual flux data. Relations that look universal, such as irrigation compensation, prove regime-specific. The contribution is not a new law of nature but a template that tells a modeller which outputs can be trusted, which need an audit, and what a discovered equation must survive before it is called one. The results concern a process model and identify the questions that measured lysimeter and eddy-covariance data can settle.

\medskip\noindent \textbf{Keywords.} equation discovery; water-balance closure; interpretable machine learning; evapotranspiration partitioning; process-model auditing\par}
\end{titlepage}

\section*{List of symbols}
\begin{table}[H]\centering\small
\caption{List of symbols, abbreviations and units used in the text.}\label{tab:notation}
\begin{tabularx}{\textwidth}{>{\raggedright\arraybackslash}p{3.3cm}Xl}
\toprule
Symbol & Meaning & Unit\\\midrule
$P_r,\ I,\ W$ & rainfall, irrigation, water input $P_r+I$ & mm\,yr$^{-1}$\\
$\mathit{ET},\ T,\ E_s,\ E_c,\ R$ & evapotranspiration and components; unattributed term, \cref{eq:R} & mm\,yr$^{-1}$\\
$D,\ \Delta S,\ Q_r$ & drainage at 1\,m, storage change, runoff & mm\,yr$^{-1}$\\
$\varepsilon$ & closure residual, \cref{eq:closure} & $-$\\
$\mathbf F,\ \sigma_k,\ \mathbf v_k$ & flux matrix, singular values and right singular vectors & --\\
$\mathbf c$ & invariant coefficients, \cref{eq:law} & $-$\\
$n,\ \alpha$ & van~Genuchten shape and scale & $-$, cm$^{-1}$\\
$\TAW$ & total available water & mm\\
$R^2_{\mathrm{cv}}$ & cross-validated coefficient of determination, \cref{eq:r2cv} & $-$\\
$\kappa$ & irrigation compensation, \cref{eq:kappa} & $-$\\
E1--E17 & numbered experiments of the Daisy ensemble & $-$\\
LOEO, LOSO & leave-one-experiment-out and leave-one-soil-out validation & $-$\\
SVD, PCA & singular-value decomposition; principal component analysis & $-$\\
GBM & gradient-boosting machine used as the flexible benchmark & $-$\\
DPF, NLF, IE & drainage fraction $D/W$; nitrogen-leaching fraction; irrigation efficiency & $-$\\\bottomrule
\end{tabularx}
\end{table}

%==========================================================
\section{Introduction}
\label{sec:intro}
Conservation of mass is the most reliable organising principle in hydrology, and it is the foundation of every water-balance model from the bucket to the Earth system model. Yet in practice the budget of a managed soil profile rarely closes exactly in the output of a process model: the fluxes that are reported separately (transpiration, soil evaporation, canopy interception) need not add to the evapotranspiration that enters the balance; change in storage is reported or implied but its size relative to the fluxes is seldom audited; and boundary fluxes at the base of a one-dimensional column are, by construction, definitions rather than measurements \citep{Hillel1998,Refsgaard2004}. At the same time, a strand of scientific machine learning has argued that conservation laws, governing equations and low-dimensional invariants can be \emph{discovered} from data, using sparse regression over libraries of candidate terms \citep{Brunton2016,Champion2019,deSilva2020}, symbolic regression \citep{Schmidt2009,Udrescu2020}, neural networks that learn invariants \citep{Iten2020,Liu2021}, and physics-informed learning \citep{Raissi2019,Karniadakis2021}. In hydrology, deep learning has delivered impressive predictive skill \citep{Kratzert2018,Kratzert2019,Shen2018}, with an ongoing debate about what such models reveal about processes \citep{Nearing2021,Reichstein2019} and about whether interpretable models should be preferred for consequential decisions \citep{Rudin2019}. Enforcing conservation analytically in learned emulators is also an active topic \citep{Beucler2021}.

Three gaps limit confidence in discovery claims in soil-water science. First, discovery methods are rarely validated against a case with a known answer before being applied to an unknown. Second, a discovered relation is typically validated within the data from which it was found, without leave-one-\emph{experiment}-out tests that ask whether the relation survives a change of regime, which is the relevant test of a ``law'' \citep{Roberts2017}. Third, negative results (where sparse laws fail to beat black-box benchmarks, or where a presumed constant is not constant) are rarely documented, so the reader cannot estimate the base rate of success \citep{Kirchner2006,Oreskes1994}.

A large designed ensemble of process-model runs, in which the true balance is known to hold up to the model's own storage term, is taken as a laboratory for method evaluation. A four-part pipeline is built (\cref{fig:pipeline}): (1) a linear-invariant search by singular-value decomposition of the flux matrix with a permutation null, used first as a positive control; (2) a unattributed-term search by auditing partitions that should sum to a total; (3) a manifold analysis of water-fate composition by compositional PCA \citep{Aitchison1986,Egozcue2003}; and (4) interpretable sparse regression over a dimensionless library with LOSO and LOEO validation against a gradient-boosting benchmark \citep{Friedman2001,Tibshirani1996,Hastie2009}. One fixed benchmark, cross-validation grouping and stopping rule are used for all targets, and every law that failed is reported.

The paper asks what a blind method can and cannot recover from a flux matrix. A linear-invariant search must first return the water balance with calibrated uncertainty, which sets the credibility of everything that follows. Accounting terms that the model's own output columns leave open are then sought, and the number of independent directions governing the partition of water among fates is counted. Compact, dimensionally sensible laws are tested against a flexible machine-learning model in \emph{unseen experiments and unseen soils}. The last aim is to separate relations that hold across regimes from those that only appear universal.

\clearpage
\newgeometry{left=0.6cm,right=0.6cm,top=0.5cm,bottom=0.5cm}
\begin{landscape}
\nolinenumbers
\begin{figure}[p]\centering
\includegraphics[width=24.6cm]{\figdir/p2/fig01_pipeline_A4.pdf}
\caption{Pipeline for interpretable discovery. The pipeline recovers the known water balance as a positive control, audits partitions, measures the dimension of the partition manifold and fits sparse laws with LOSO/LOEO validation against gradient boosting. Controls (dashed purple boxes) are shown at the bottom and outputs at the top; purple hexagons are processing blocks, the pink stippled box is the data store, peach boxes are outputs and dotted lines are flows of information. The 149 runs excluded as numerical failures are analysed in \cref{sec:failed}.}
\label{fig:pipeline}
\end{figure}
\end{landscape}
\restoregeometry


\paragraph{Study design and levels of evidence.} The unit of observation is the annual-mean water balance of one simulated run, and the runs of an experiment are correlated because they share a design; for this reason validation withholds whole experiments and whole soils. Four levels of evidence are kept apart. Computational reproducibility, the regeneration of every reported number from archived code and data, is established. Predictive validity is assessed by leave-one-experiment-out and leave-one-soil-out scoring against a flexible benchmark, with a nested evaluation of the selection step. Mechanistic validity is limited to the process model, so that the unattributed flux is a property of the model's accounting. Field transferability is not tested, and the conclusions are stated as questions that measured data can settle.

%==========================================================
\section{Materials and methods}
\label{sec:methods}
\subsection{Data}
The Daisy \citep{Hansen2012,Holbak2026} virtual-experiment ensemble is used, as described in the companion paper: 10\,190 runs in 17 experiments (E0--E16), ten soils (S01--S10, van~Genuchten $n=1.23$--2.68, \cref{tab:soilsP2}), 20-year Taastrup weather; 149 runs failed and are excluded (\cref{sec:failed}), and 10\,041 valid runs are analysed. Each run contributes 20-year mean annual rainfall $P_r$, irrigation $I$, evapotranspiration $\mathit{ET}$, transpiration $T$, soil evaporation $E_s$, canopy evaporation $E_c$, drainage at 1\,m $D$, and (derived) crop yield, nitrogen leaching and irrigation efficiency. Experiments differ in structure: single-soil (coarse loam) lattices (E8, E11, E12, E14, E15), all-soil factorials and lattices (E9, E16) and one-factor series (E1--E4). The distributions of the key variables are in \cref{tab:desc}.

\begin{table}[H]\centering\small
\caption{Descriptive statistics of the 10\,041 valid runs (mm\,yr$^{-1}$; yield in Mg\,ha$^{-1}$ not shown).}
\label{tab:desc}
\begin{tabular}{lrrrrr}
\toprule
Variable & Mean & SD & 25\,\% & Median & 75\,\%\\\midrule
Rain $P_r$ & 417 & 143 & 297 & 464 & 464\\
Irrigation $I$ & 180 & 191 & 0 & 178 & 268\\
Water input $W=P_r+I$ & 597 & 251 & 464 & 654 & 722\\
Evapotranspiration $\mathit{ET}$ & 546 & 196 & 386 & 611 & 671\\
Drainage $D$ (1\,m) & 60 & 106 & 3 & 34 & 78\\
Unattributed $R$ (\cref{eq:R}) & 164 & 97 & 103 & 129 & 223\\\bottomrule
\end{tabular}
\end{table}

\begin{table}[H]\centering\small
\caption{Soil parameters used as predictors (effective van~Genuchten values from the project documentation; $\TAW$ in mm to 1\,m).}
\label{tab:soilsP2}
\begin{tabular}{lllrr}
\toprule
ID & Texture & $\alpha$ (cm$^{-1}$) & $n$ & $\TAW$\\\midrule
S01 & Coarse loam & 0.0961 & 1.353 & 130\\
S02 & Coarse sand & 0.0536 & 1.600 & 105\\
S03 & Sand & 0.1450 & 2.680 & 60\\
S04 & Loamy sand & 0.1240 & 2.280 & 75\\
S05 & Sandy loam & 0.0750 & 1.890 & 95\\
S06 & Loam & 0.0360 & 1.560 & 150\\
S07 & Silt loam & 0.0200 & 1.410 & 220\\
S08 & Sandy clay loam & 0.0590 & 1.480 & 160\\
S09 & Clay loam & 0.0190 & 1.310 & 175\\
S10 & Silty clay loam & 0.0100 & 1.230 & 195\\\bottomrule
\end{tabular}
\end{table}

\subsection{Failed runs and why they are excluded}
\label{sec:failed}
Of the 10\,190 runs, 149 (1.5\,\%) are flagged as failed in the ensemble and are excluded from every analysis. Because exclusion that is not random can bias a discovery study, the pattern of failure was examined before the runs were dropped, and the argument is set out in four parts.

\paragraph{Failures are concentrated in three soils and in specific experiments.} Of the 149 failures, 129 (87\,\%) occur in the coarse sand S02, 13 in the sand S03 and 7 in the silty clay loam S10; the other seven soils have none (\cref{fig:failed}a). In S02, 24.9\,\% of the 519 runs fail. In S03 all 13 failures belong to the initial-moisture experiment E9, and in S10 all 7 belong to the nitrogen--water and dose experiments (E13, E1). The three soils lie at the extremes or in the anomalous part of the van~Genuchten range ($n=1.60$ for S02, $n=2.68$ for S03, the steepest retention curve, and $n=1.23$ for S10, the flattest), which is where the unsaturated conductivity function changes most rapidly with suction.

\paragraph{Within a failing soil, failure follows the size of the water pulse.} In S02, the fraction of failed runs rises with the irrigation depth applied per event in experiment E10: none of the 20 runs at 1 and 2\,mm fails, 10\,\% at 3\,mm, 50\,\% at 8--10\,mm and all of the 30 runs at 15, 20 and 30\,mm (\cref{fig:failed}b). It also depends on the rain scale: 69\,\% of the S02 runs at a scale of 1.5 fail (22 of 32) and 12\,\% at 0.5 (5 of 42), against none of the 90 runs at intermediate scales 0.6--1.4 (\cref{fig:failed}c). Runs that share the same weather, the same soil and the same code, and differ only in how much water arrives at once, therefore fail or succeed in a systematic order. This is the behaviour of a state-dependent numerical difficulty and not of a corrupted input, which would hit runs irrespective of their forcing.

\paragraph{The outputs of failed runs are physically impossible.} In valid runs the annual drainage at 1\,m is never below $-17.5$\,mm (median in valid S02 runs, 134\,mm). In failed runs it is negative in 65\,\% of cases and falls to $-882$\,mm (\cref{fig:failed}d), and the ratio of evapotranspiration to water input ($P+I$) ranges from 0.54 to 4.34, with a median of 1.09, so that in some runs more water leaves the profile than entered it. A net upward flux of hundreds of millimetres through the lower boundary of a free-draining column, or evapotranspiration several times the water input, cannot be produced by the physics of the model; it is the signature of a solver that has not converged, in which a mass error is absorbed in the boundary or storage terms.

\paragraph{The alternatives are less consistent with the pattern.} A fault in the weather or in the run definition would not depend on soil and on dose in this ordered way; a failure from over-wetting in the field, such as ponding or runoff, is a physical process that the model is built to represent and would not give negative drainage; and a random job failure would spread evenly over soils. Numerical non-convergence of the Richards-equation solver in a stiff regime (a very permeable, quickly draining column that receives a large input, or a column with an extreme retention curve) accounts for the concentration in soils, the order with water input and the impossible outputs together.

\paragraph{Treatment of the excluded runs.} The failed runs are interpreted as numerical failures of the soil-water solver under stiff conditions, not as information about soil hydrology, and they are excluded for that reason; the interpretation follows from the pattern of failure described above. Two consequences follow. The exclusion removes a quarter of the runs of S02, preferentially in the wettest and driest forcing and at the largest doses, so results for this soil, and any law fitted to it, apply to the part of the factor space in which the solver succeeded. The other nine soils are essentially unaffected (at most 2.5\,\% of runs in S03), and the conclusions that rest on across-soil comparisons were tested for sensitivity to the exclusion by re-analysis. With S02 removed (9\,651 runs) the three-term ET law is re-estimated with coefficients $0.652$ for $\ln W$, $-0.297$ for $n$ and $0.029$ for $\ln(1+I)$, against $0.646$, $-0.299$ and $0.029$ for the full ensemble, with LOEO $R^2=0.897$ and LOSO $R^2=0.904$ (full ensemble 0.898 and 0.904), and the unattributed term $R$ is 30\,\% of ET (163\,mm\,yr$^{-1}$; full ensemble 30\,\% and 164\,mm\,yr$^{-1}$). With S02 and S03 both removed (9\,145 runs) the scores are 0.893 and 0.898 and $R$ is 30\,\% (165\,mm\,yr$^{-1}$). The exclusion therefore does not drive the across-soil results.

\begin{figure}[H]\centering
\includegraphics[width=\textwidth]{\figdir/p2/figS_failed.pdf}
\caption{Pattern of the 149 failed runs. (a) Failed runs per soil, as a percentage of the 519 runs per soil (counts above bars). (b) Failed fraction in S02 against irrigation depth per event in E10. (c) Failed fraction in S02 against rain scale. (d) Annual drainage at 1\,m for valid and failed runs, normalised as densities; negative drainage of several hundred millimetres occurs only in failed runs.}
\label{fig:failed}
\end{figure}

\subsection{Governing balance and the flux matrix}
For each run the mean annual water balance is
\begin{equation}
P_r+I \;=\; \mathit{ET}+D+\Delta S+Q_r,\qquad \mathit{ET}=T+E_s+E_c+R,
\label{eq:wb2}
\end{equation}
where $\Delta S$ is storage change, $Q_r$ runoff, and $R$ is defined as the remainder of ET after the three reported components,
\begin{equation}
R \;=\; \mathit{ET}-T-E_s-E_c .
\label{eq:R}
\end{equation}
Equation~\eqref{eq:R} defines $R$ without ascribing a process. Let $\mathbf F\in\mathbb R^{N\times q}$ be the matrix whose rows are runs and whose columns are the fluxes (rain, irrigation, ET, $T$, $E_s$, $E_c$, $D$). If a linear invariant $\sum_j c_jF_{ij}=0$ holds for all runs, then $\mathbf c$ lies in the null space of $\mathbf F$. Define $W=P_r+I$ and the closure residual
\begin{equation}
\varepsilon \;=\; \frac{W-\mathit{ET}-D}{W}\;\approx\;\frac{\Delta S+Q_r}{W}.
\label{eq:closure}
\end{equation}

\subsection{Linear-invariant search by SVD with a permutation null}
The singular-value decomposition $\mathbf F=\mathbf U\boldsymbol\Sigma\mathbf V^{\top}$ with singular values $\sigma_1\ge\dots\ge\sigma_q$ identifies near-null directions as the right singular vectors $\mathbf v_k$ with small $\sigma_k$. Each column is first divided by its root-mean-square value so that no flux dominates through its units; $\mathbf v_k$ is mapped back to mm units and normalised so that the ET coefficient equals $-1$. Because singular values depend on the scale of the columns, significance is judged against a \emph{column-permutation null}: each column of $\mathbf F$ is permuted independently across runs (destroying cross-column dependence but preserving each marginal), and the smallest singular value under the null is recorded; an invariant is accepted only if $\sigma_{\min}$ lies far below the null value. Uncertainty of the coefficients is estimated by a nonparametric bootstrap over runs (resampling 10\,041 rows; 95\,\% percentile intervals). The test is repeated (a) on the four-flux subsystem $(P_r,I,\mathit{ET},D)$, which is the set that appears in the balance; (b) on the seven-column matrix; and (c) excluding E8 and E9, whose initial states are deliberately away from equilibrium. It is emphasised that this is a \emph{positive control}: the expected answer is $\mathbf c=(1,1,-1,-1)$ up to scale and the pipeline is judged by how closely it recovers it, and by whether it fails in a diagnosable way.

Any decomposition that should sum to a total is audited by computing the residual. The audit covers (a) \cref{eq:R}, reporting $R$, its share $R/\mathit{ET}$ by experiment, and its dependence on drivers by regression of $\ln R$ on candidate drivers (irrigation, soil, mulch capacity, temperature, evaporation factor); and (b) the storage closure $\varepsilon$ of \cref{eq:closure} by experiment. The audit is descriptive: it identifies where the budget does not close, and tests candidate attributions by whether $R$ responds to the drivers that would be expected under each.

The partition of water input into fates (transpiration, soil evaporation, canopy evaporation, unattributed $R$, drainage) is compositional: components are positive and sum to a constant, so Euclidean statistics are inappropriate \citep{Aitchison1986,Egozcue2003}. Writing the composition as $\mathbf x=(x_1,\dots,x_5)$ with $x_i=\text{flux}_i/\sum_j\text{flux}_j$, the centred log-ratio (clr) transform is applied
\begin{equation}
\mathrm{clr}(\mathbf x)_i=\ln x_i-\frac1{5}\sum_{j=1}^5\ln x_j ,
\label{eq:clr}
\end{equation}
and PCA \citep{Wold1987} to the transformed data. Drainage is zero or negative in 14\,\% of runs, so zeros are replaced by a floor $\delta$ (in mm\,yr$^{-1}$) before the closure to unit sum; because the results depend strongly on $\delta$, results are reported for $\delta\in\{10^{-3},0.1,1,5,10\}$. Because log-transformations are inherently sensitive to near-zero values, the sensitivity to the floor is mapped across four orders of magnitude and not fixed at a single value, so that the stability of the drainage axis can be judged independently of the floor selected. The number of components needed to explain the variance is a data-driven estimate of the dimension of the partition manifold.

For a target $y$ (log of a flux), a library $\boldsymbol\Theta$ is built of dimensionless or log-transformed candidates (\cref{tab:lib}). The model is
\begin{equation}
\ln y \;=\; \beta_0+\sum_{k\in\mathcal S}\beta_k\,\theta_k+\epsilon,\qquad \mathcal S\subseteq\{1,\dots,p\},
\label{eq:sparse}
\end{equation}
with the active set $\mathcal S$ chosen by forward selection: at each iteration the candidate that maximises the leave-one-experiment-out cross-validated $R^2$ is added, and the search stops when the improvement is below 0.003. The cross-validated $R^2$ is
\begin{equation}
R^2_{\mathrm{cv}}=1-\frac{\sum_i(\ln y_i-\widehat{\ln y}_i^{(-g(i))})^2}{\sum_i(\ln y_i-\overline{\ln y})^2},
\label{eq:r2cv}
\end{equation}
where $\widehat{\ln y}_i^{(-g(i))}$ is the prediction for run $i$ from a model fitted with all runs of its group $g(i)$ held out. Groups are experiments (LOEO, 17 folds) or soils (LOSO, 10 folds). LOEO is the demanding test: it asks whether a law fitted on 16 experimental designs predicts the 17th, which differs in which factors vary and in what range. The root-mean-square error in mm and the mean absolute percentage error after back-transformation are also reported. The benchmark is a histogram-based gradient-boosting regressor \citep{Friedman2001,Pedregosa2011} with 250 iterations on the whole library, validated under the same groupings. Caveat: because forward selection uses the LOEO score, the reported LOEO $R^2$ of the selected model is mildly optimistic (selection bias); LOSO values, which were not used for selection, provide a partly independent check, and coefficient stability across folds is reported.

\begin{table}[H]\centering\small
\caption{Candidate library $\boldsymbol\Theta$ (log-transformed or dimensionless).}
\label{tab:lib}
\begin{tabular}{lll}
\toprule
Symbol & Definition & Role\\\midrule
$\ln P_r$ & log rainfall & climate input\\
$\ln(1+I)$ & log irrigation & management input\\
$\ln W$ & log water input $P_r+I$ & total supply\\
$\ln E_p$ & log soil-evaporation factor & evaporative demand\\
$\Delta T$ & temperature shift (\textdegree C) & warming\\
$n$ & van~Genuchten shape parameter & pore-size distribution\\
$\ln\TAW$ & log total available water & storage capacity\\
$\ln\alpha$ & log van~Genuchten scaling & air-entry scale\\
mulch & VFF$\times$CAP (mulch lattice) & residue layer\\
$I/W$ & irrigation fraction & managed share\\
$\ln(W/\TAW)$ & dimensionless supply-to-storage & wetness index\\\bottomrule
\end{tabular}
\end{table}

\label{sec:svdtheory}
Suppose the true model is $\mathbf F\mathbf c=\boldsymbol\eta$ with $\boldsymbol\eta$ a small run-wise residual (here $-\Delta S-Q_r$ for $\mathbf c=(1,1,-1,-1)$). The SVD gives the best rank-$(q-1)$ approximation of $\mathbf F$ (Eckart--Young), and the right singular vector $\mathbf v_q$ minimises $\|\mathbf F\mathbf v\|_2$ over unit vectors; the corresponding singular value is $\sigma_q=\|\mathbf F\mathbf v_q\|_2$. Because the columns are all measured with error (here, the model's own discretisation and storage noise), $\mathbf v_q$ is a total-least-squares solution, in which errors in all variables are treated symmetrically; this is preferable to regressing one flux on the others, whose coefficients would be attenuated and depend on the choice of dependent variable. When $\eta$ is uncorrelated with the columns, $\sigma_q^2\approx\|\boldsymbol\eta\|^2$, so $\sigma_q/\sqrt N$ is the root-mean-square residual per run in the units of the scaled columns. Before the decomposition each column is divided by its root-mean-square value (so that no flux dominates by its units), and the singular vector is mapped back to mm units afterwards; the singular values are therefore dimensionless, and for the four-flux subsystem $\sigma_q/\sqrt N=5.9/\sqrt{10041}=0.059$, i.e.\ a residual of about 6\,\% of a typical column scale. A caveat is that the matrix is not mean-centred: centring would remove the constant direction, but a conservation law $\sum c_jF_j=0$ has no intercept, so the uncentred matrix is the correct object.

Permuting each column independently preserves marginal distributions and destroys all linear relations between columns, so the smallest singular value under permutation, $\sigma_q^{\mathrm{perm}}$, measures how small $\sigma_q$ would be if the columns were unrelated but had the same scales. A ratio $\sigma_q/\sigma_q^{\mathrm{perm}}=5.9/33.3=0.18$ for the four-flux matrix shows that the null direction reflects genuine dependence. Because fluxes are positively skewed and share a common experiment structure (for example rainfall is identical across many runs), the permutation null is conservative in one respect (it breaks the experimental clustering) and liberal in another (it does not respect the clustering when permuting). Runs are therefore also bootstrapped and intervals reported.

\begin{enumerate}[nosep,label=\arabic*.]
\item Initialise $\mathcal S=\emptyset$, best score $s^\ast=-\infty$.
\item For each candidate $k\notin\mathcal S$, fit \cref{eq:sparse} on the training groups and compute $R^2_{\mathrm{cv}}$ (\cref{eq:r2cv}) for $\mathcal S\cup\{k\}$ with 17 LOEO folds.
\item Let $k^\ast$ be the maximiser; if $R^2_{\mathrm{cv}}(\mathcal S\cup\{k^\ast\})>s^\ast+0.003$, add $k^\ast$ and update $s^\ast$; otherwise stop.
\item Refit on all data to report coefficients; evaluate LOSO using the same terms.
\end{enumerate}
\Cref{tab:path} shows the path for the ET law: $\ln W$ alone gives $R^2_{\mathrm{cv}}=0.801$, adding $n$ raises it to 0.895 and adding $\ln(1+I)$ to 0.898. The third term improves the score by only 0.003, at the stopping threshold, so a two-term law $\ln\mathit{ET}=f(\ln W,n)$ is almost as good and is the safer summary.

The compensation coefficient $\kappa$ measures the share of a rainfall change that is offset by irrigation. For rainfall change $\Delta P_r$ with other factors fixed,
\begin{equation}
\kappa \;=\; -\frac{\Delta I}{\Delta P_r}.
\label{eq:kappa}
\end{equation}
In single-soil lattices $\mathrm d W/\mathrm dP_r=1-\kappa$ is estimated by ordinary least squares; in E16, paired differences are formed of rain scale 0.5 and 1.5 across the 32 combinations of other factors within each soil, and bootstrap the pooled estimate (500 resamples).

All analyses are scripted in Python \citep{Harris2020,Virtanen2020,Pedregosa2011,Hunter2007}; scripts and intermediate results are in the project repository (\datalink, \codelink) and mapped in Appendix~\ref{app:repro}. Model evaluation follows the recommendations of \citet{Moriasi2007} and \citet{Bennett2013}; $R^2$ is reported on the log scale unless stated. No test data were used for model selection other than through the declared cross-validation.

%==========================================================
\section{Results}
\label{sec:results}
\subsection{The pipeline recovers the water balance (positive control)}
The singular spectrum of the seven-flux matrix falls steeply, with a smallest singular value of 5.5 compared with about 33 under column permutation (\cref{fig:invariants}a), and the corresponding right singular vector places weight on rain, irrigation, ET and drainage with opposite signs for inputs and outputs (\cref{fig:invariants}b). For the four-flux subsystem the null direction, scaled to $c_{\mathit{ET}}=-1$, is
\begin{equation}
\mathbf c=(c_{P_r},c_I,c_{\mathit{ET}},c_D)=(1.037,\;0.957,\;-1,\;-1.036),
\label{eq:law}
\end{equation}
with bootstrap 95\,\% intervals 1.031--1.045 for rain, 0.940--0.974 for irrigation and $-1.052$ to $-1.023$ for drainage, and a smallest singular value of 5.9 against 33.3 under permutation. The pipeline therefore recovers the conservation law $P_r+I-\mathit{ET}-D\approx0$ with every coefficient within 5\,\% of the exact value, and the permutation null shows that the null direction is not an artefact of scale. \emph{However, the bootstrap intervals do not contain the exact values}: rain $\pm$ 1 lies below the rain interval, and irrigation lies above its interval. The deviation reflects the nonzero storage term and other unreported fluxes: excluding E8 and E9, whose initial moisture is deliberately unbalanced, gives $(1.021,\,0.969,\,-1,\,-1.019)$ with a smallest singular value of 3.0, closer to the exact balance but still not exact.

\begin{figure}[H]\centering
\includegraphics[width=\textwidth]{\figdir/p2/fig02_invariants.pdf}
\caption{Linear invariants. (a) Singular values of the seven-flux matrix for the real data and for column-permuted data. (b) Coefficients of the smallest right singular vector (scaled so $c_{\mathit{ET}}=-1$) with bootstrap intervals. (c) Storage residual $\Delta S/W$ by experiment: E8 and E9 do not close. (d) Unattributed share of ET, $R/\mathit{ET}$, by experiment.}
\label{fig:invariants}
\end{figure}

The result illustrates both the promise and the discipline needed for law discovery. The method finds the law that is there, without being told which variables are inputs and which are outputs; but its coefficients should not be over-interpreted, because the third decimal of a coefficient reflects what the data do \emph{not} close, not what the law is.

Fitting the four-flux null vector separately within each experiment (\cref{tab:svdexp}) tests whether the pooled vector \cref{eq:law} hides heterogeneity. In 11 of the 14 experiments with irrigation all coefficients are within about 5\,\% of the exact values $(1,1,-1,-1)$ (E11 is at 5.1\,\% for drainage; E16, with the extreme corners of the factorial, gives 0.901 for irrigation and $-0.919$ for drainage, E6 gives 0.961, 0.927 and $-0.869$; and E12 gives drainage $-1.116$ are the exceptions). The residual scale $\sigma_q/\sqrt N$ is 0.001--0.017 in the well-closing experiments and 0.081 in E16, 0.030 in E6, 0.166 in E8 and 0.219 in E9. E0 (no irrigation variation within the 10 baseline runs), E8 and E9 (rainfed) have a constant zero irrigation column and the invariant is formed from three fluxes; for E8 the drainage coefficient is $+15.5$, a sign of an ill-conditioned decomposition in an experiment in which storage dominates the balance, which is a cautionary example of how an unconstrained search can return a numerically valid but physically meaningless vector. The method therefore identifies its own failure by the magnitude of $\sigma_q/\sqrt N$ and by physical-plausibility checks on the coefficients, which is recommended as a rule.

\begin{table}[H]\centering\small
\caption{Four-flux null vector fitted within each experiment (scaled to $c_{\mathit{ET}}=-1$); ``--'' marks a flux with no variation in the experiment (irrigation in E0, E8, E9). $\hat\sigma=\sigma_q/\sqrt N$ is the residual scale in units of the column root-mean-square.}
\label{tab:svdexp}
\begin{tabular}{lrrrrr}
\toprule
Experiment & Runs & $c_{P_r}$ & $c_{I}$ & $c_{D}$ & $\hat\sigma$\\\midrule
E0 & 10 & 1.025 & -- & $-$1.041 & 0.004\\
E1 & 189 & 1.009 & 0.969 & $-$0.979 & 0.017\\
E2 & 197 & 1.022 & 0.983 & $-$1.027 & 0.004\\
E3 & 200 & 1.016 & 0.995 & $-$1.012 & 0.004\\
E4 & 96 & 1.006 & 0.989 & $-$0.983 & 0.003\\
E5 & 300 & 1.016 & 0.993 & $-$1.027 & 0.002\\
E6 & 193 & 0.961 & 0.927 & $-$0.869 & 0.030\\
E7 & 200 & 1.008 & 0.992 & $-$0.990 & 0.003\\
E8 & 1000 & 1.241 & -- & 15.496 & 0.166\\
E9 & 987 & 1.182 & -- & $-$1.249 & 0.219\\
E10 & 954 & 1.026 & 0.972 & $-$1.038 & 0.010\\
E11 & 1000 & 0.998 & 1.024 & $-$1.051 & 0.004\\
E12 & 1000 & 1.016 & 1.005 & $-$1.116 & 0.005\\
E13 & 975 & 1.016 & 0.984 & $-$1.008 & 0.004\\
E14 & 1000 & 0.982 & 1.035 & $-$0.974 & 0.001\\
E15 & 1000 & 1.009 & 0.966 & $-$0.971 & 0.003\\
E16 & 740 & 1.018 & 0.901 & $-$0.919 & 0.081\\
\bottomrule
\end{tabular}
\end{table}

The closure residual $\varepsilon$ (\cref{eq:closure}) has a median absolute value of 0.9\,\% of water input over all valid runs and 0.65\,\% excluding E8 and E9 (95th percentile 2.3\,\%), so for 15 of the 17 experiments the mass balance holds to within a few percent. In E8 and E9, which vary initial moisture, the median absolute residual is 16\,\% of $W$, with the residual reaching $\pm0.3$--0.6 in the tails (\cref{fig:invariants}c). This is expected, because these experiments deliberately impose a non-equilibrium initial storage that is released over the 20-year window, but it matters for discovery: a conservation law fitted without storage will appear to be violated in exactly the experiments where hydraulic memory is strongest. A global law should therefore include storage explicitly or be fitted only on near-equilibrium experiments, and the two cases give different coefficients (\cref{eq:law} and the text following it).

\subsection{A unattributed term in the ET partition}
The ET components that Daisy reports do not add up to ET. The remainder $R$ (\cref{eq:R}) averages 164\,mm\,yr$^{-1}$, with an interquartile range of 103--223\,mm\,yr$^{-1}$, and equals 30\,\% of ET (median 29\,\%; range 2--88\,\%). Its share varies between experiments from 14\,\% (E6, E8) to 41\,\% (E7, E14; \cref{fig:invariants}d, \cref{tab:R}). The dependence of $R$ on management is informative (\cref{fig:hidden}). In E4 (methods), $R$ is 105\,mm\,yr$^{-1}$ for every subsurface-drip configuration (5--60\,cm), 183 for surface and 166 for overhead irrigation, so that \emph{where} water is applied matters while the drip depth does not (\cref{fig:hidden}a). In the rainfed E9 lattice, $R=50.9+0.121\,P_r$ (\cref{fig:hidden}b), i.e.\ about 12\,\% of rainfall, and in E11 (coarse loam)
\begin{equation}
\ln R=5.75-0.355\ln P+0.863\ln E_p+0.136\,\Delta T,\qquad R^2=0.979,
\label{eq:Rlaw}
\end{equation}
where $P$ is the rain scale, $E_p$ the soil-evaporation factor and $\Delta T$ the temperature shift, so that $R$ scales almost linearly with the soil-evaporation factor, decreases with rain scale (elasticity $-0.36$) and increases with warming. The mulch experiments E7 and E14 have the highest median shares (41\,\%) and the cross-experiment sparse model for $\ln R$ selects $\ln(1+I)$, $n$, the mulch capacity and $\Delta T$.

\begin{figure}[H]\centering
\includegraphics[width=\textwidth]{\figdir/p2/fig03_hidden_term.pdf}
\caption{The unattributed term $R$. (a) $R$ and canopy evaporation in E4 by irrigation method. (b) $R$ against annual rain in E9 (rainfed, light points), with the linear fit and the ten E0 baselines.}
\label{fig:hidden}
\end{figure}

\begin{table}[H]\centering\small
\caption{Unattributed term $R$ by experiment. median $R$ (mm\,yr$^{-1}$), median share of ET, and median storage residual $\Delta S/W$.}
\label{tab:R}
\begin{tabular}{lrrrr}
\toprule
Experiment & Runs & $R$ (mm) & $R/\mathit{ET}$ & $\Delta S/W$\\\midrule
E0 Baselines & 10 & 107 & 0.24 & $-0.021$\\
E1 Doses & 189 & 118 & 0.21 & $-0.001$\\
E4 Methods & 96 & 103 & 0.19 & $-0.003$\\
E5 SHP dynamics & 300 & 162 & 0.29 & $-0.007$\\
E6 Mediterranean & 193 & 70 & 0.14 & $-0.012$\\
E7 Mulching & 200 & 226 & 0.41 & $-0.006$\\
E8 Rain intensity$\times$moisture & 1000 & 39 & 0.14 & $-0.133$\\
E9 Soil moisture$\times$rain & 987 & 81 & 0.37 & $-0.014$\\
E10 Amount, timing, soil & 954 & 118 & 0.23 & $-0.011$\\
E11 Temperature$\times E_p\times$rain & 1000 & 296 & 0.38 & $-0.009$\\
E12 Sensor, trigger, rain & 1000 & 218 & 0.32 & $-0.008$\\
E13 N$\times$water & 975 & 125 & 0.28 & $-0.006$\\
E14 Mulch$\times$rain & 1000 & 262 & 0.41 & $0.001$\\
E15 Method, depth, amount & 1000 & 103 & 0.17 & 0.003\\
E16 Edge cases & 740 & 157 & 0.30 & $-0.010$\\\bottomrule
\end{tabular}
\end{table}

What is $R$? The data do not identify it, but they constrain candidates. A process that is independent of drip depth, larger under surface and overhead application, proportional to the soil-evaporation factor, increased under mulch and weakly increasing with warming is consistent with evaporation from the \emph{surface layer} (mulch or litter and ponded water) that is not reported in the soil-evaporation or canopy-evaporation columns. That reading is a hypothesis; the model code has not been inspected to confirm which Daisy state variables sum into the model's ET output, and a software-level audit of the output definitions would be required to turn the hypothesis into a statement. The practical consequence is nonetheless clear: any analysis that treats $T+E_s+E_c$ as ET (as many reduced-form ET-partitioning analyses do) omits about 30\,\% of the flux, and the partition fractions reported for the ensemble depend on how $R$ is allocated.

A sparse model for $R$ across all experiments has limited out-of-sample skill. The forward-selected model $\ln R=3.35+0.194\ln(1+I)+0.492\,n+0.145\,\mathrm{cap}+0.238\,\Delta T$ has an in-sample $R^2=0.61$ on the log scale but a leave-one-experiment-out error of 0.51 in $\ln R$, equal to a level-scale $R^2$ of only 0.30 (RMSE 81\,mm) versus 0.50 (RMSE 69\,mm) for gradient boosting. Within a single experiment, \cref{eq:Rlaw} is excellent; across experiments it is not, which implies that $R$ depends on interactions among irrigation method, mulch state and demand that the library does not contain.

\subsection{A dominant drainage axis, with its variance share quantified by bootstrap}
The log-ratio transform requires strictly positive components, and drainage is zero or negative (net capillary rise) in 1\,429 of the 10\,041 runs (14\,\%) and below 1\,mm\,yr$^{-1}$ in 1\,947. The treatment of these zeros is therefore not a detail. The first analysis replaced all components below $10^{-3}$\,mm\,yr$^{-1}$ by that floor, which gave explained variances of 88.6\,\%, 7.3\,\% and 3.3\,\% for the first three components. Because $\ln(10^{-3})$ lies far from the logs of the other components, the runs with zero drainage then dominate the first axis. Repeating the analysis with floors of 0.1, 1, 5 and 10\,mm\,yr$^{-1}$ gives a first-component share of 74.3, 62.7, 56.5 and 52.3\,\% (\cref{fig:manifold}a, \cref{tab:clipsens}). The result that a single drainage axis dominates is robust, since the loading of drainage on PC1 stays between $+0.82$ and $+0.88$ for every floor and also when only runs with drainage above 5 or 20\,mm are used (PC1 share 57\,\% and 52\,\%), but its \emph{magnitude} is a function of the data-handling choice and the figure of 88.6\,\% should not be quoted as a property of the system. The number of components needed to explain 90\,\% of the variance is two to three in all cases (three components carry 92--98\,\% for floors from $10^{-3}$ to 10\,mm, with 97.3\,\% at the 1\,mm floor adopted for the loadings, \cref{fig:manifold}b,c).

With the 1\,mm floor, the first component (62.7\,\%) is a drainage-versus-evaporation axis: drainage loads $+0.84$, soil evaporation $-0.37$, unattributed $R$ $-0.34$, transpiration $-0.20$ and canopy evaporation $+0.06$. Its scores correlate with the drainage fraction $D/W$ at Spearman $\rho=0.98$, with $\ln(W/\TAW)$ at 0.48 and with $n$ at 0.45. The second component (24.8\,\%) separates canopy evaporation ($+0.77$) from soil evaporation ($-0.56$), and the third (9.8\,\%) opposes $R$ ($+0.83$) to soil evaporation ($-0.43$) and transpiration ($-0.26$). The choice of the 1\,mm floor was made \emph{after} seeing the sensitivity and is justified only as a value close to the resolution of annual flux accounting; the full sensitivity is reported so that readers can judge. The practical conclusion is that the drainage fraction is, to a first approximation, the leading summary of water partitioning (its score correlation is 0.98), and that a second axis, the canopy-versus-soil evaporation split, is needed for a complete description.

\begin{figure}[H]\centering
\includegraphics[width=\textwidth]{\figdir/p2/fig04_manifold.pdf}
\caption{Partition manifold and its sensitivity to zero handling. (a) Variance explained by the first three clr-PCA components for five zero-replacement floors. (b) Scores on the first two components (1\,mm floor) coloured by $\ln(W/\TAW)$. (c) Loadings of the five fates on the first two components.}
\label{fig:manifold}
\end{figure}

\begin{table}[H]\centering\small
\caption{Sensitivity of the clr-PCA to the zero-replacement floor (variance explained, \%).}
\label{tab:clipsens}
\begin{tabular}{lrrrr}
\toprule
Floor (mm\,yr$^{-1}$) & PC1 & PC2 & PC3 & PC4\\\midrule
$10^{-3}$ & 88.6 & 7.3 & 3.3 & 0.8\\
0.1 & 74.3 & 17.3 & 6.6 & 1.7\\
1 & 62.7 & 24.8 & 9.8 & 2.6\\
5 & 56.5 & 25.1 & 14.8 & 3.6\\
10 & 52.3 & 27.1 & 16.8 & 3.8\\\bottomrule
\end{tabular}
\end{table}

Averaged over the 15 near-equilibrium experiments, the five fates (\cref{fig:fatesoil}a) show the expected gradient: drainage takes 34.5\,\% of water input in sand (S03), 28.5\,\% in loamy sand and 5.4\,\% in coarse loam (S01), 8.5\,\% in silty clay loam, while transpiration rises from 19\,\% in sand to 31--33\,\% in the loams. The unattributed term, in contrast, is remarkably constant across soils: $R$ takes 22--26\,\% of $W$ in nine of ten soils (31\,\% in coarse loam, which has the mulch and temperature lattices) and its share of ET is 23--39\,\% (\cref{fig:fatesoil}c), the largest in sand, where there is least else to evaporate. Since $R$ is almost soil-independent in share of $W$, it behaves as a surface-driven term rather than a soil-hydraulic one, consistent with the surface-layer hypothesis. The stacked fractions sum to 0.96--1.01 of $W$ (the remainder being storage change and runoff). Panel~b is a reminder that coarse loam contributes 4\,419 of the near-equilibrium runs while each other soil contributes 290--419.

\begin{figure}[H]\centering
\includegraphics[width=\textwidth]{\figdir/p2/fig07_fates_by_soil.pdf}
\caption{Water fates by soil in the 15 near-equilibrium experiments (E8, E9 excluded). (a) Mean fate shares of water input $W$, soils ordered by decreasing $n$. (b) Number of valid runs per soil. (c) Median share of the unattributed term in ET.}
\label{fig:fatesoil}
\end{figure}

\subsection{A compact evapotranspiration law with out-of-sample skill}
Forward selection over the library selects three terms for ET (\cref{tab:laws}): $\ln W$, $n$ and $\ln(1+I)$, giving
\begin{equation}
\ln\mathit{ET}=2.503+0.646\,\ln W-0.299\,n+0.029\,\ln(1+I),
\label{eq:ETlaw}
\end{equation}
with ordinary-least-squares standard errors $0.027$, $0.0051$, $0.0045$ and $0.0011$ (which understate uncertainty because runs are not independent) and in-sample $R^2=0.908$. Leave-one-experiment-out $R^2=0.898$ and leave-one-soil-out $R^2=0.904$, against 0.876--0.879 (LOEO) and 0.881--0.894 (LOSO) for gradient boosting in two independent runs, so the three-term law matches or exceeds the benchmark under both groupings (\cref{fig:ETlaw}a,b). The RMSE in leave-one-experiment-out prediction is 68.8\,mm\,yr$^{-1}$ (MAPE 9.2\,\%). The fitted coefficients are stable across the 17 folds (\cref{fig:ETlaw}c; coefficient means 2.51, 0.645, $-0.298$, 0.029 with standard deviations of 0.077, 0.017, 0.017, 0.0023).

Equation~\eqref{eq:ETlaw} has a straightforward physical reading. The exponent on $W$ (0.646) means that ET rises sublinearly with water supply, the signature of an energy- or storage-limited regime also encoded in Budyko-type curves \citep{Budyko1974,Milly1994,Zhang2001}: a 10\,\% increase in supply raises ET by 6.5\,\%, and the remainder goes to drainage, runoff and storage. The coefficient on $n$ (-0.299) implies a factor $\exp(-0.299\times1.45)=0.65$ in ET between the most uniform (S03, $n=2.68$) and the most graded (S10, $n=1.23$) soil at the same $W$: soils that drain freely retain less water for evaporation, in agreement with the hydraulic-memory results of the companion paper. The small positive coefficient on $\ln(1+I)$ (0.029) indicates that, at equal total input, a larger irrigated share raises ET slightly, plausibly through the wetting of the surface that increases $R$ and $E_s$. The law is \emph{not} a fundamental conservation law, nor a model of mechanism; it is a compact summary of how the model's ET responds across 17 designs.

\begin{figure}[H]\centering
\includegraphics[width=\textwidth]{\figdir/p2/fig05_ET_law.pdf}
\caption{Compact ET law. (a) Leave-one-experiment-out predictions of the three-term law versus Daisy ET. (b) Gradient-boosting benchmark. (c) Fitted coefficients across the 17 LOEO folds. (d) LOEO MAPE by experiment. (e) The law at $I=200$\,mm for four soils (grey dashed: ET $=W$). (f) Drainage derived from closure, $W-\widehat{\mathit{ET}}$, versus Daisy drainage ($R^2=0.69$, RMSE 64\,mm; E8 and E9 excluded).}
\label{fig:ETlaw}
\end{figure}

The 17 experiments were resampled with replacement (2\,000 cluster-bootstrap resamples) and the cross-validated predictions of the law were compared with those of a gradient-boosting model that uses the same three inputs ($\ln W$, $n$, $\ln(1+I)$; LOEO $R^2=0.866$). The paired difference in $R^2$ is $+0.032$ with a 95\,\% interval of 0.006--0.080, so the law is better than a flexible learner given the same information in this resampling, with the caveat that the law's terms were selected with the same folds. With the \emph{full} library the benchmark reaches 0.876--0.879, and the two are regarded as comparable rather than claiming superiority.

The error is not uniform (\cref{tab:byexp}). Mean absolute percentage errors are 2.6--5.8\,\% in most one- and two-factor experiments, but 16.6\,\% in E16 (edge cases: extreme combinations), 16.7\,\% in E8 and 27.4\,\% in E9, the two experiments that violate storage closure, and 11.7\,\% in E5 (time-varying soil hydraulic properties, where a single static $n$ is wrong by construction). The law is therefore valid where the water balance is near equilibrium and soil hydraulics are static; it should not be extrapolated to initial-state perturbations or to dynamic soil structure.

\begin{table}[H]\centering\small
\caption{Leave-one-experiment-out performance of \cref{eq:ETlaw} by experiment: MAPE (\%) and mean error (predicted minus Daisy, mm\,yr$^{-1}$).}
\label{tab:byexp}
\begin{tabular}{lrrr}
\toprule
Experiment & Runs & MAPE (\%) & Mean error (mm)\\\midrule
E0 Baselines & 10 & 5.8 & $-25.2$\\
E1 Doses & 189 & 3.8 & 9.8\\
E2 Trigger pressure & 197 & 3.0 & $-6.6$\\
E3 Sensor depth & 200 & 2.6 & $-8.0$\\
E4 Methods & 96 & 3.6 & 9.4\\
E5 SHP dynamics & 300 & 11.7 & $-63.8$\\
E6 Mediterranean & 193 & 8.5 & 41.5\\
E7 Mulching & 200 & 4.7 & 10.6\\
E8 Rain intensity & 1000 & 16.7 & $-11.3$\\
E9 Soil moisture & 987 & 27.4 & 12.9\\
E10 Amount, timing & 954 & 3.3 & 1.3\\
E11 Temperature, $E_p$, rain & 1000 & 9.0 & $-62.6$\\
E12 Sensor, trigger, rain & 1000 & 3.3 & $-15.2$\\
E13 N, water & 975 & 3.5 & $-3.7$\\
E14 Mulch, rain & 1000 & 4.0 & 7.5\\
E15 Method, depth, amount & 1000 & 5.0 & 29.0\\
E16 Edge cases & 740 & 16.6 & 34.5\\\bottomrule
\end{tabular}
\end{table}

\Cref{tab:closure2} gives the closure residual $\varepsilon$ for every experiment. In 15 experiments the median absolute residual is below 2.1\,\% of $W$ and the 95th percentile below 4\,\% (E6, the Mediterranean-pollution scenarios, is the exception at 10.7\,\%); in E8 and E9 the median absolute residual is about 16\,\% and the 95th percentile exceeds 100\,\% of $W$. For E8 the mean (signed) residual is negative, $-13\,\%$: the profile is gaining water, that is, the run begins drier than the equilibrium state and refills. The law \cref{eq:ETlaw} was fitted on water inputs from 124 to 2\,762\,mm\,yr$^{-1}$ (ET of 95--1\,600\,mm\,yr$^{-1}$), $n$ from 1.23 to 2.68, in a temperate climate, and it should be applied only within this domain and with near-equilibrium storage.

\begin{table}[H]\centering\small
\caption{Closure residual $\varepsilon=(W-\mathit{ET}-D)/W$ by experiment (per cent of water input): median signed, median absolute and 95th percentile of absolute residual.}
\label{tab:closure2}
\begin{tabular}{lrrrr}
\toprule
Experiment & Runs & Median $\varepsilon$ & Median $|\varepsilon|$ & 95th pct.\ $|\varepsilon|$\\\midrule
E0 & 10 & -2.06 & 2.06 & 2.9\\
E1 & 189 & -0.12 & 0.33 & 1.2\\
E2 & 197 & -0.87 & 0.87 & 2.3\\
E3 & 200 & -0.87 & 0.87 & 2.0\\
E4 & 96 & -0.33 & 0.34 & 1.2\\
E5 & 300 & -0.74 & 0.74 & 1.9\\
E6 & 193 & -1.19 & 1.57 & 10.7\\
E7 & 200 & -0.57 & 0.57 & 1.2\\
E8 & 1000 & -13.35 & 15.76 & 105.8\\
E9 & 987 & -1.40 & 16.87 & 108.5\\
E10 & 954 & -1.08 & 1.17 & 2.5\\
E11 & 1000 & -0.93 & 0.94 & 1.8\\
E12 & 1000 & -0.82 & 0.83 & 1.8\\
E13 & 975 & -0.64 & 0.64 & 2.3\\
E14 & 1000 & 0.09 & 0.40 & 1.5\\
E15 & 1000 & 0.34 & 0.34 & 0.4\\
E16 & 740 & -1.01 & 1.22 & 3.8\\
\bottomrule
\end{tabular}
\end{table}

The ensemble is heavily imbalanced: 5\,519 of the 10\,041 runs (55\,\%) are coarse loam (S01) because the large lattices E8, E11, E12, E14 and E15 use only that soil, so soil-level behaviour in the other nine soils rests on about 400--520 runs each. The LOEO residuals of \cref{eq:ETlaw} by soil (\cref{tab:soilres}) show that the law over-predicts ET in the fine soils (mean error $+22$, $+20$ and $+31$\,mm in S08, S09 and S10, where positive means over-prediction relative to Daisy, since errors are predicted minus Daisy) and has the largest percentage errors in sand (15.1\,\%) and loamy sand (13.7\,\%), the soils for which $n$ is most extreme. A single shape parameter therefore does not remove all soil dependence, and the MAPE of about 7.5--15\,\% by soil is consistent with the overall 9.2\,\%.

\begin{table}[H]\centering\small
\caption{LOEO residuals of the ET law by soil (predicted minus Daisy, mm\,yr$^{-1}$). MAPE: mean absolute percentage error; RMSE in mm\,yr$^{-1}$.}
\label{tab:soilres}
\begin{tabular}{llrrrr}
\toprule
Soil & Texture & Runs & Mean error & MAPE (\%) & RMSE\\\midrule
S01 & Coarse loam & 5519 & -10.7 & 7.5 & 58.9\\
S02 & Coarse sand & 390 & 15.8 & 8.6 & 69.8\\
S03 & Sand & 506 & -13.6 & 15.1 & 63.8\\
S04 & Loamy sand & 519 & -4.8 & 13.7 & 57.9\\
S05 & Sandy loam & 519 & -3.8 & 10.9 & 56.1\\
S06 & Loam & 519 & -0.2 & 9.9 & 77.5\\
S07 & Silt loam & 519 & 0.4 & 10.6 & 106.2\\
S08 & Sandy clay loam & 519 & 22.2 & 10.2 & 57.8\\
S09 & Clay loam & 519 & 19.7 & 9.9 & 92.8\\
S10 & Silty clay loam & 512 & 31.2 & 11.6 & 107.4\\
\bottomrule
\end{tabular}
\end{table}

\begin{table}[H]\centering\small
\caption{Forward-selection path for the ET law. LOEO $R^2_{\mathrm{cv}}$ after each added term; and for other targets (first terms).}
\label{tab:path}
\begin{tabular}{llr}
\toprule
Target & Term added & $R^2_{\mathrm{cv}}$ (LOEO)\\\midrule
$\ln\mathit{ET}$ & 1 ($\ln W$) & 0.801\\
 & 2 ($n$) & 0.895\\
 & 3 ($\ln(1+I)$) & 0.898\\
$\ln T$ & 1 ($\ln W$) & 0.581\\
 & 2 ($n$) & 0.641\\
$\ln E_s$ & 1 ($n$) & 0.070\\
 & 2 (mulch) & 0.237\\
 & 3 ($\ln P_r$) & 0.308\\
$\ln E_c$ & 1 ($\ln W$) & 0.641\\
 & 2 ($\ln P_r$) & 0.651\\
 & 3 ($I/W$) & 0.668\\
$\ln$ yield & 1 ($\ln W$) & 0.531\\
 & 2 ($\ln(1+I)$) & 0.549\\
 & 3 ($\ln P_r$) & 0.599\\
 & 4 (mulch) & 0.604\\
$\ln$ DPF & 1 ($\ln(W/\TAW)$) & 0.193\\
 & 2 ($\ln\alpha$) & 0.260\\
 & 3 ($I/W$) & 0.422\\
 & 4 ($\Delta T$) & 0.433\\\bottomrule
\end{tabular}
\end{table}

If E8 and E9, which break closure, are removed from both fitting and validation (15 experiments, 8\,054 runs), the same three terms give LOEO $R^2=0.882$, LOSO $R^2=0.881$, RMSE 65.9\,mm and MAPE 5.8\,\%, so the law is more accurate (relative error 5.8\,\% versus 9.2\,\%) in the domain for which it is intended. The marginal contribution of the terms is clearer in this subset: $\ln W$ alone gives LOEO $R^2=0.600$ (RMSE 138\,mm), adding $n$ raises it to 0.857 (RMSE 74\,mm) and adding $\ln(1+I)$ to 0.882. A gradient-boosting model restricted to the same three inputs reaches LOEO $R^2=0.778$, below the law, so the advantage of the compact form is not that it uses more information but that its functional form (power law in $W$, exponential in $n$) extrapolates across experiments better than piecewise-constant trees, which cannot extrapolate beyond the range of the training experiments. (This comparison uses three inputs and 15 folds and is not directly comparable with the full-library benchmark of \cref{tab:laws}.)

\Cref{tab:example} applies \cref{eq:ETlaw} to three combinations of water input and irrigation in three soils. At equal input, ET is about 35\,\% lower in sand than in silty clay loam ($263$ versus $406$\,mm at $W=400$\,mm), because water not held in the root zone drains; the implied drainage plus storage change $W-\mathit{ET}$ is 137\,mm in sand and negative-to-small in the fine soils, in line with the memory results. The example is illustrative; its uncertainty is the RMSE of 66--69\,mm.

\begin{table}[H]\centering\small
\caption{ET (mm\,yr$^{-1}$) predicted by \cref{eq:ETlaw} for three water inputs $W$ with irrigation $I$.}
\label{tab:example}
\begin{tabular}{lrrr}
\toprule
 & $W=400$, $I=0$ & $W=650$, $I=200$ & $W=1000$, $I=300$\\\midrule
S03 sand ($n=2.68$) & 263 & 420 & 561\\
S01 coarse loam ($n=1.353$) & 391 & 624 & 834\\
S10 silty clay loam ($n=1.23$) & 406 & 648 & 865\\\bottomrule
\end{tabular}
\end{table}

\label{sec:failures}
\Cref{tab:laws} collects all sparse laws that were fitted. Transpiration is moderately predictable by a two-term law ($\ln T=-1.329+1.127\ln W-0.501\,n$; LOEO $R^2=0.641$, equal to gradient boosting 0.642). Canopy evaporation is predicted at LOEO 0.668 (gradient boosting 0.691) by $\ln E_c=-13.3-0.470\ln W+3.157\ln P_r+4.240\,I/W$, whose large coefficients reflect the narrow range of $E_c$ and should not be interpreted mechanistically. Crop yield is predicted at LOEO 0.604 (gradient boosting 0.650). In contrast, soil evaporation is poorly predicted by sparse terms ($\ln E_s=3.271-0.770\,n-0.380\,\mathrm{mulch}+0.462\ln P_r$; LOEO 0.308 versus 0.547 for gradient boosting; LOSO 0.348 versus 0.476), indicating that interactions that the library does not contain control $E_s$. The drainage fraction (DPF $=D/W$), the natural target given the PCA result, is the worst: the best sparse model with exogenous predictors has a LOEO $R^2$ of 0.43 and a LOSO $R^2$ of 0.22, against 0.50 (LOEO) for gradient boosting, and allowing partition predictors in the library raises gradient boosting to 0.59 but does not change the sparse law. The shortfall is a negative result for compact drainage laws.

\begin{table}[H]\centering\small
\caption{Sparse laws and benchmark. cross-validated $R^2$ on the log scale. Selected terms in order of selection.}
\label{tab:laws}
\footnotesize\setlength{\tabcolsep}{4pt}
\begin{tabular}{llrrrr}
\toprule
Target & Selected terms & Sparse LOEO & Sparse LOSO & GBM LOEO & GBM LOSO\\\midrule
$\ln\mathit{ET}$ & $\ln W$, $n$, $\ln(1+I)$ & 0.898 & 0.904 & 0.876--0.879 & 0.881--0.894\\
$\ln T$ & $\ln W$, $n$ & 0.641 & 0.652 & 0.642 & 0.659\\
$\ln E_s$ & $n$, mulch, $\ln P_r$ & 0.308 & 0.348 & 0.547 & 0.476\\
$\ln E_c$ & $\ln W$, $\ln P_r$, $I/W$ & 0.668 & 0.669 & 0.691 & 0.806\\
$\ln$ yield & $\ln W$, $\ln(1+I)$, $\ln P_r$, mulch & 0.604 & 0.617 & 0.650 & 0.743\\
$\ln$ DPF & $\ln(W/\TAW)$, $\ln\alpha$, $I/W$, $\Delta T$ & 0.433 & 0.217 & 0.498 & --\\
$\ln R$ & $\ln(1+I)$, $n$, cap, $\Delta T$ & 0.41$^{\dagger}$ & 0.43$^{\dagger}$ & 0.43--0.47$^{\dagger}$ & --\\\bottomrule
\end{tabular}

\smallskip\noindent\footnotesize$^{\dagger}$ $R^2$ on the log scale for $R$; level-scale $R^2$ is 0.30 (sparse) and 0.50--0.51 (gradient boosting).
\end{table}

Because the balance holds to within a few percent in 15 of 17 experiments, drainage can be derived by closure as $\hat D=W-\widehat{\mathit{ET}}$, which converts a hard prediction problem (drainage) into one with a good compact law (ET). On the 8\,054 runs outside E8 and E9 this yields $R^2=0.69$ and RMSE 64\,mm against Daisy drainage (\cref{fig:ETlaw}f), and negative predicted drainage in 4.8\,\% of runs. The accuracy is limited by the relative size of the quantities and is not uniform (\cref{tab:dbins}): the mean absolute error is smallest, 21--28\,mm, for Daisy drainage between 10 and 200\,mm (relative error 19--66\,\%), but the closure-derived value overpredicts by 32\,mm on average for drainage below 10\,mm (1\,944 runs) and \emph{underpredicts} by 69\,mm and 308\,mm for drainage of 200--500 and above 500\,mm. The ET law is smooth in $W$ and cannot represent the large drainage events, so closure-derived drainage is a useful consistency check for moderate drainage and not a predictor of extremes.

\begin{table}[H]\centering\small
\caption{Closure-derived drainage $W-\widehat{\mathit{ET}}$ (LOEO ET law) against Daisy drainage by drainage class (E8, E9 excluded). Bias = predicted minus Daisy.}
\label{tab:dbins}
\begin{tabular}{lrrr}
\toprule
Daisy drainage (mm\,yr$^{-1}$) & Runs & Bias (mm) & Mean absolute error (mm)\\\midrule
$<10$ & 1944 & $+31.7$ & 39.7\\
10--50 & 2631 & $+8.2$ & 21.0\\
50--100 & 1736 & $-2.2$ & 25.6\\
100--200 & 1284 & $-13.7$ & 28.3\\
200--500 & 377 & $-68.8$ & 104.0\\
$>500$ & 82 & $-308.1$ & 318.6\\\bottomrule
\end{tabular}
\end{table}

\subsection{Irrigation compensation is regime-specific}
A natural candidate invariant is the compensation of rainfall changes by irrigation, $\kappa$ (\cref{eq:kappa}). It is not universal (\cref{fig:comp}, \cref{tab:kappa}). In single-soil lattices the slope of $W$ on rain is $0.58\pm0.02$ in E11 ($\kappa=0.42$) and 0.33 in E12 and E14 ($\kappa=0.67$), whereas in E16 the pooled $\kappa$ across the ten soils is 0.159 (bootstrap 95\,\% interval 0.134--0.181) and the per-soil values range from $-0.014$ (sand) to 0.303 (coarse loam). Within E16, $\kappa$ falls from 0.193 to 0.123 as $E_p$ is raised from 0.3 to 1.2, from 0.188 to 0.129 with warming from $+2$ to $+6$\,\textdegree C, and from 0.211 to 0.107 with trigger suction from 50 to 2000\,cm, but rises from 0.079 to 0.241 with irrigation dose from 1.2 to 30\,mm. Compensation therefore depends on the scheduling rule and soil, and a single value of $\kappa$ should not be used in water-allocation arguments.

\begin{figure}[H]\centering
\includegraphics[width=\textwidth]{\figdir/p2/fig06_compensation.pdf}
\caption{Irrigation compensation $\kappa$ is not universal. (a) Per-soil $\kappa$ in E16 (paired $P=0.5\to1.5$) with the pooled estimate and bootstrap interval. (b) Slope of water input on rain in single-soil lattices (95\,\% intervals). (c) Mean $\kappa$ by factor level in E16.}
\label{fig:comp}
\end{figure}

\begin{table}[H]\centering\small
\caption{Compensation $\kappa=-\Delta I/\Delta P_r$ and the response of ET and drainage to rain, per soil in E16.}
\label{tab:kappa}
\begin{tabular}{lrrrr}
\toprule
Soil & $\kappa$ & $\Delta\mathit{ET}/\Delta P_r$ & $\Delta D/\Delta P_r$ & Pairs\\\midrule
S01 & 0.303 & 0.512 & 0.186 & 32\\
S02 & 0.164 & 0.346 & 0.495 & 8\\
S03 & $-0.014$ & 0.473 & 0.550 & 32\\
S04 & 0.084 & 0.473 & 0.453 & 32\\
S05 & 0.102 & 0.551 & 0.358 & 32\\
S06 & 0.108 & 0.619 & 0.281 & 32\\
S07 & 0.216 & 0.542 & 0.246 & 32\\
S08 & 0.090 & 0.613 & 0.308 & 32\\
S09 & 0.238 & 0.515 & 0.253 & 32\\
S10 & 0.302 & 0.477 & 0.237 & 32\\\bottomrule
\end{tabular}
\end{table}

%==========================================================
\section{Discussion}
\label{sec:disc}
\subsection{Contributions of a validated discovery protocol and a quantified accounting gap}
The paper contributes a discovery protocol in which every stage is tested against a known answer, and a first application that exposes an unclosed term in a widely used process model. The pipeline is blind to variable roles; it recovers the water balance as the dominant invariant, with a permutation null (smallest singular value 5.9 against 33.3 for permuted data), and this positive control is what licenses the interpretation of every other result. Applied to the evapotranspiration partition, the same pipeline quantifies a term that the reported components do not carry, $R\approx164$\,mm\,yr$^{-1}$ or about 30\,\% of ET, and a storage non-closure in two experiments (median 16\,\%). A compact three-term law then predicts ET across withheld experiments and withheld soils (LOEO $R^2=0.898$, LOSO 0.904) as accurately as gradient boosting (0.879 and 0.894). Quantifying a closure gap has a long record as a productive step in flux science: the energy-balance closure problem at eddy-covariance sites, in which turbulent fluxes typically fall 10--30\,\% short of available energy, was characterised across a network and then drove advances in measurement and in the interpretation of flux data \citep{Wilson2002,Foken2008}. The unattributed term $R$ has the same function for model-output accounting: it converts an implicit assumption (that reported components are the total) into a measured quantity that can be traced to a pathway. Stating the recovered invariant for what it is, the known mass balance returned by a blind method, is the credential that makes the new findings credible in a field where discovery claims are often overstated \citep{Rudin2019,Nearing2021}.

Evapotranspiration partitioning is a central topic, from global studies that infer transpiration fractions \citep{Jasechko2013,Schlesinger2014} to field studies of mulch effects \citep{BalwinderSingh2011}. If a model's reported ET components omit a third of the flux, then a partition fraction computed as $T/(T+E_s+E_c)$ is larger than $T/\mathit{ET}$ by a factor $1/(1-0.30)\approx1.4$ on average, an error comparable with the differences between management treatments that are routinely reported. The size of $R$ also helps to explain why the closure-derived and component-based estimates of soil evaporation differ. For users of Daisy and similar models, the practical recommendation is to compute ET as the model's own total and to report $R$ alongside the components, and to document, in the model output specification, exactly which processes enter the total \citep{Holbak2026}.

\subsection{A map of predictability across the flux budget}
The ET law works because ET is the large, smooth, sub-linear function of supply that the water-balance literature has described for decades, and because the model's ET is energy- and storage-limited \citep{Budyko1974,Milly1994,Zhang2001}. The drainage fraction is the small residual of two large numbers, it is strongly nonlinear in the wetness of the profile, and it depends on thresholds (field capacity, trigger) that a log-linear library cannot represent \citep{Rodriguez2004}; gradient boosting captures some of this and still only reaches $R^2\approx0.5$ in LOEO. This locates the information boundary: the content that drainage needs lies beyond the annual-flux library (e.g.\ intra-annual timing, event intensity, antecedent state). The companion paper shows that antecedent state is indeed decisive for drainage in coarse soils.

The three-term law is competitive with gradient boosting under LOEO and LOSO, which supports the view that accuracy does not always require opacity \citep{Rudin2019}. The law does not extend to the regimes in which the physics it omits is first-order (initial-state experiments, dynamic hydraulic properties), and the gradient-boosting benchmark fails there too. This agreement is itself informative: it shows that the boundary is set by the information in the library and not by the model class, which tells a modeller which targets can be emulated by compact laws and which need timing and state variables.

Equation discovery in soil hydrology has so far targeted the governing partial differential equation, notably the Richards equation recovered from soil-moisture time series by group-sparse regression \citep{Song2022,Song2025}, and deep-learning studies have imposed water-balance closure as a constraint on forecasts \citep{Li2024}. The present paper is complementary and differs in its object and its emphasis: it works with annual flux budgets, tests whether a pipeline recovers a known balance before interpreting any other result, audits the closure of the partition that supplies its targets, and validates by withheld regimes. The present approach shares with SINDy-style discovery \citep{Brunton2016} the use of sparse regression over a library, but differs in two respects: it targets steady annual relations rather than dynamics, and it places the validation burden on leave-one-regime-out testing \citep{Roberts2017}. It shares with physics-informed learning \citep{Raissi2019,Karniadakis2021} the emphasis on conservation, but imposes it as a \emph{diagnostic} (closure residual) rather than as a constraint \citep{Beucler2021}; applying it as a constraint would have hidden $R$, which is the finding most relevant to model users. Related work on learning conserved quantities \citep{Liu2021,Iten2020} would be a natural extension to daily dynamics.

Emulators are used to make process-model knowledge available in fast decision tools \citep{Razavi2012,Asher2015}. The results here indicate which emulator targets are safe: ET and transpiration are predictable from compact laws; drainage and soil evaporation are not, and should be emulated with care or by closure. The companion software paper applies these lessons to the annual meta-model embedded in SWAFL Desktop and quantifies its rank-faithful (rather than absolute) accuracy.

A software-level audit would resolve the identity of $R$ directly. Four questions, in order of cost: (i) Which state variables are summed into the output called ``evapotranspiration'' in the log file used here, and which into the three reported components? (ii) Is evaporation from surface ponds, mulch and litter reported in a separate column that was not exported in the ensemble? (iii) Does the canopy-evaporation column represent interception only, excluding evaporation of water retained on residue? (iv) Is the storage term in the model's own water balance consistent with the exported depth and boundary definitions (for example, is the 1\,m drainage column the flux at the bottom of the root zone or the soil column)? If the answer to (ii) is yes, the unattributed term would be explained and the diagnostic here would have served its purpose as an audit.

Random $k$-fold cross-validation of the same models would have produced much higher scores, because runs inside an experiment are close neighbours in factor space. To illustrate the effect on the present data: a model that predicts $\ln\mathit{ET}$ from $\ln W$ alone has LOEO $R^2$ of 0.80 on the full data and 0.60 when E8 and E9 are removed. Leave-one-experiment-out removes entire designs, so the model must extrapolate to unseen combinations of factor ranges, which is exactly what a ``law'' is supposed to do \citep{Roberts2017}. The leave-one-soil-out scores answer a second question: would the law hold for a soil not used to fit it? That the two scores agree for ET (0.898 and 0.904) is evidence that the law is not driven by a single experiment or soil, but it is not evidence of validity in other climates. A nested evaluation, in which the forward selection was repeated inside each training set of 16 experiments and scored only on the held-out experiment, gives a pooled $R^2=0.895$, against 0.898 for the single-level score, and selects the same three terms in 15 of 17 outer folds (the two-term set $\ln W$, $n$ in the other two); the optimism from selecting on the evaluation score is therefore small for this law.

The sub-linear exponent $\partial\ln\mathit{ET}/\partial\ln W=0.646$ can be compared with the elasticity of Budyko-type curves. For a Budyko curve $\mathit{ET}/P=F(\Phi)$ with aridity index $\Phi=\mathit{PET}/P$, the elasticity of ET to $P$ lies between 0 (energy limit) and 1 (water limit) and is about 0.5--0.8 at intermediate aridity \citep{Milly1994,Zhang2001,Budyko1974}. The fitted value is within that range. Unlike the Budyko approach, which treats catchment properties via a single parameter, the present law includes a soil term ($n$): the fitted coefficient implies that a unit increase in $n$ reduces ET by 26\,\% ($1-e^{-0.299}$) at fixed $W$, which can be read as the dependence of the effective Budyko parameter on soil hydraulic uniformity. The law is not claimed to be a Budyko curve; the mapping would require PET as a variable, and PET is constant within the weather series used.


A sequential protocol follows from these results and is offered as a general template for discovery studies in hydrology. The first element is a positive control: the pipeline is applied to a system whose governing balance is known, and it must recover that balance with a null-referenced uncertainty, as the permutation test and bootstrap do here for the water balance \citep{Efron1993,Good2005}. The second is a closure audit of every partition used as a target, since a decomposition that omits a term will propagate that omission into every law fitted to it, as the unattributed term $R$ propagates into the transpiration and soil-evaporation fractions. The third is regime-wise validation, in which whole experiments, soils or climates are withheld, because random folds inside a designed ensemble measure interpolation between neighbouring factor combinations and not transfer to a new regime \citep{Roberts2017}. The fourth is a comparison with a flexible benchmark, reported with the cases in which the benchmark is better, so that the interpretability premium can be read directly \citep{Rudin2019,Friedman2001}. The fifth is a stated domain of validity for each law, expressed in the experiments and soils for which the leave-one-out score exceeds a stated threshold.

The unattributed term $R$ also defines a concrete verification programme for the process model. Three comparisons would identify it. First, an ablation in which individual surface and canopy evaporation pathways are switched off in turn would show which pathway removes $R$ from the balance. Second, a tracer-style accounting that records every term added to the evaporative flux in the model log, with its state variable of origin, would establish whether $R$ is a reporting convention or a physical pathway. Third, a comparison across structurally different models (Richards-based, capacity-based and radiation-driven) on the same designed ensemble would show whether a non-closing partition is a feature of one model family or a property of how evapotranspiration is partitioned in general \citep{Simunek2008,Jones2003,Keating2003,Steduto2009}. The experimental data that would constrain the physical interpretation are lysimeter and eddy-covariance records in which evaporation of intercepted and ponded water is measured separately from transpiration \citep{Good2015,Jasechko2013,Schlesinger2014}.

Beyond hydrology, the same sequence of controls is relevant to other settings in which a regression of a conserved quantity on its drivers is read as a discovered relation. Sparse regression over a library is attractive because the result is an equation that can be examined, dimensionally checked and compared with theory \citep{Brunton2016,Udrescu2020,Champion2019}. Its reliability, however, depends on the design of the data. In an ensemble built by crossing factors, several drivers are correlated by construction, and the selection path of a forward algorithm reflects that correlation; the leave-one-experiment-out score then serves two roles, as selection criterion and as validation, which is why the LOSO scores, which are less affected by selection, are reported alongside it. The practical consequence is that the coefficients of a discovered law should be reported with the validation regime in which they were obtained, and that a law should be described as discovered only for the regimes in which it is validated.

Finally, the results bear on the design of the emulators and decision tools that draw on process-model ensembles. A tool that reports evapotranspiration components should state whether they sum to the total; one that reports drainage should state whether it is simulated or derived by closure, and the evidence here shows that closure-derived drainage is accurate in sign and rank order but under-estimates large drainage. The companion software article applies these requirements by labelling each equation with its provenance and by testing the emulator against held-out simulations \citep{Razavi2012,Asher2015,Wilkinson2016}.

\subsection{Domain of validity of the discovered laws}
Each law is reported with the regimes in which it was validated, and each boundary defines the next test.
\begin{enumerate}[nosep]
\item \emph{Process-model domain.} The results describe Daisy output. $R$ is a property of the model's accounting; its counterpart in nature is measured in lysimeters and eddy-covariance records that separate intercepted and ponded water from transpiration.
\item \emph{Failed runs.} The 149 excluded runs (87\,\% in S02) form a pattern that identifies them as solver failures (\cref{sec:failed}); re-running with tighter tolerances is the confirming step, and S02 is under-represented at large doses and extreme rain scales.
\item \emph{Model and weather.} One model and one climate; the laws are conditional on both, and the leave-one-regime-out protocol is designed to be repeated on further models and climates.
\item \emph{Selection.} Forward selection used the LOEO score; the nested evaluation, with selection repeated inside each training set, gives $R^2=0.895$ against 0.898 for the ET law, so the selection optimism is small. LOSO values are less affected by selection.
\item \emph{Dependence of runs.} Coefficient standard errors assume independent runs; cross-fold variation is the reported guide to uncertainty.
\item \emph{Soil covariates.} Ten soils; $n$ carries several correlated soil properties.
\item \emph{Identification of $R$.} Its surface-layer origin is inferred from its driver dependence, and the four-question audit of the model code in this section settles it.
\item \emph{Aggregation.} Twenty-year means average over seasonality and event structure; the companion work addresses antecedent state.
\item \emph{Benchmark reproducibility.} The gradient-boosting benchmark is run with a fixed random seed. Across unseeded repeats the LOSO score varied by up to 0.013, so differences below about 0.02 between a law and the benchmark are not interpreted.
\end{enumerate}

(1) Begin with a positive control with a known answer. (2) Use a permutation null for any decomposition-based invariant. (3) Validate across \emph{regimes} (experiments, soils, climates), not only across random folds. (4) Audit partitions for closure before using them as targets. (5) Compare with a flexible benchmark and report where the benchmark wins. (6) State a domain of validity for each law. (7) Publish failures.

%==========================================================
\section{Conclusions}
A blind linear-invariant search recovers the water balance from a flux matrix with every coefficient within 5\,\% of the exact value, calibrated against a permutation null, and flags two accounting gaps: an unattributed evapotranspiration term averaging 164\,mm\,yr$^{-1}$ (30\,\% of ET) that depends on irrigation method, evaporation factor and mulch, and a storage non-closure of 16\,\% in the two initial-moisture experiments. Compositional PCA shows that a single drainage--evaporation axis dominates water partitioning, with a variance share (52--89\,\%) that depends strongly on how zero drainage is handled. A three-term law, $\ln\mathit{ET}=2.503+0.646\ln W-0.299\,n+0.029\ln(1+I)$, matches gradient boosting in leave-one-experiment-out ($R^2=0.898$) and leave-one-soil-out ($R^2=0.904$) validation, while sparse laws for drainage fraction, soil evaporation and the unattributed term do not, and irrigation compensation is not a universal constant. The conclusions are limited to a single process model and weather series, and the discovered ``law'' is the known balance together with a diagnosed accounting gap.

\section*{Data, code and software availability}
\codeavail

\section*{Author contributions}
N.H.: conceptualisation, analysis, software, visualisation, writing.

\section*{Funding}
No funding was received for this work.

\section*{Competing interests}
The author declares no competing interests.

\bibliography{refs}

\appendix
\section{Supplementary information}
The analysis was exploratory and was not pre-registered. The main analytical decisions were: the four controls (permutation null, bootstrap, LOSO/LOEO, gradient-boosting benchmark); the forward-selection stopping threshold (0.003); log-targets; the exclusion of E8 and E9 as a sensitivity analysis for the invariant search and for closure-derived drainage; and the treatment of $R$ as a remainder. The candidate library was assembled by the author with knowledge of the ensemble and of the companion hydraulic-memory results (for example, the choice to include $n$, $\ln\alpha$ and $\ln W/\TAW$), which increases researcher degrees of freedom; readers should treat LOEO and LOSO scores as indicative rather than as confirmatory error rates.

The seven-flux search also returns weights on transpiration, soil evaporation and canopy evaporation in its smallest direction (rain 0.896, irrigation 0.834, ET $-1$, transpiration 0.367, soil evaporation 0.121, canopy evaporation $-0.122$, drainage $-0.933$), which is not a clean invariant: because $T+E_s+E_c+R=\mathit{ET}$ and $R$ is not a column, the seven-column matrix has an additional near-dependence that mixes components. This is why the four-flux system is the clean positive control.

\begin{table}[H]\centering\small
\caption{Columns of the ensemble used in this paper.}
\begin{tabularx}{\textwidth}{llX}
\toprule
Column & Unit & Definition\\\midrule
\texttt{experiment}, \texttt{treatment}, \texttt{run\_id} & -- & experiment code, treatment string encoding factor levels, unique run identifier\\
\texttt{soil} & -- & S01--S10\\
\texttt{rain\_mm}, \texttt{irrigation\_mm} & mm\,yr$^{-1}$ & 20-year mean annual rainfall and irrigation\\
\texttt{ET\_mm} & mm\,yr$^{-1}$ & model evapotranspiration (total)\\
\texttt{transpiration\_mm}, \texttt{soil\_evap\_mm}, \texttt{canopy\_evap\_mm} & mm\,yr$^{-1}$ & reported components of ET\\
\texttt{drainage\_1m\_mm} & mm\,yr$^{-1}$ & flux across the 1\,m boundary (negative = net rise)\\
\texttt{water\_stress\_d} & d\,yr$^{-1}$ & days of water stress\\
\texttt{yield\_Mg\_ha} & Mg\,ha$^{-1}$ & grain/biomass yield\\
\texttt{N\_fertiliser\_kg\_ha}, \texttt{N\_leached\_kg\_ha} & kg\,ha$^{-1}$ & nitrogen applied and leached\\
\texttt{IE\_kg\_m3}, \texttt{DPF}, \texttt{NLF} & -- & irrigation efficiency, drainage fraction $D/W$, N-leaching fraction (definitional; excluded as targets)\\
\texttt{P}, \texttt{dT}, \texttt{Ep}, \texttt{trig}, \texttt{z}, \texttt{amt\_mm}, \texttt{vff}, \texttt{cap}, \texttt{Nrate} & various & factor levels (rain scale, warming, evaporation factor, trigger, sensor depth, dose, bare fraction, mulch capacity, N rate)\\
\bottomrule
\end{tabularx}
\end{table}

DPF$=D/W$ is exact by construction and its ``prediction'' is therefore the same task as predicting $D$; NLF and IE are ratios of other columns. They are not used as independent evidence in any result.

\section{Software views of the water-balance partition}
The partition analysed in this paper is displayed in the SWAFL Desktop interface, documented in the software \manuallink, which can be reviewed or downloaded as a PDF. \Cref{fig:swshots} shows the field water balance and the layers that expose the partition of evapotranspiration into transpiration, soil evaporation and canopy evaporation (captured at 4000$\times$2250 px; files in \texttt{figures/shots/} of the repository).
\begin{figure}[H]\centering
\includegraphics[width=0.92\textwidth]{\figdir/shots/layer_n1.png}\\[6pt]
\includegraphics[width=0.92\textwidth]{\figdir/shots/layer_n2.png}\\[6pt]
\includegraphics[width=0.92\textwidth]{\figdir/shots/dash_results.png}
\caption{SWAFL Desktop views of the quantities analysed in this paper: water-balance layers (upper and middle) and the field water balance with the results panel (lower).}
\label{fig:swshots}
\end{figure}

\section{Reproducibility map}\label{app:repro}
\begin{table}[H]\centering\footnotesize
\caption{Result-to-script map.}
\begin{tabularx}{\textwidth}{XX}
\toprule
Result & Script / output\\\midrule
Ensemble table and run counts & \texttt{make\_data.py} $\to$ \texttt{data/daisy\_ensemble.csv.gz}\\
SVD, permutation null, closure (\cref{eq:law}, \cref{fig:invariants}) & \texttt{p2\_a.py}, \texttt{p2\_laws.py} $\to$ \texttt{p2\_a.json}, \texttt{p2\_laws.json}\\
Unattributed term $R$ (\cref{eq:Rlaw}, \cref{fig:hidden}) & \texttt{p2\_R.py} $\to$ \texttt{p2\_R.json}\\
clr-PCA (\cref{fig:manifold}, \cref{tab:clipsens}) & \texttt{p2\_a.py}, \texttt{p2\_figs.py}, \texttt{p2\_fig4.py} $\to$ \texttt{p2\_pca.json}, \texttt{p2\_pca\_sens.json}\\
Sparse laws (\cref{tab:laws}, \cref{eq:ETlaw}) & \texttt{p2\_et.py}, \texttt{p2\_dpf.py}, \texttt{p2\_figs2.py}, \texttt{p2\_nested.py} $\to$ \texttt{p2\_et.json}, \texttt{p2\_dpf.json}, \texttt{p2\_etlaw.json}, \texttt{p2\_nested.json}; exclusion sensitivity \texttt{p2\_s02\_excl.py} $\to$ \texttt{p2\_s02\_excl.json}\\
Compensation (\cref{fig:comp}) & \texttt{p2\_laws2.py} $\to$ \texttt{p2\_laws2.json}\\
Figures & \texttt{p2\_figs.py}, \texttt{p2\_figs2.py}, \texttt{p2\_fig4.py}, \texttt{p2\_fig6.py}, \texttt{p2\_fig7.py}, \texttt{p2\_figfail.py}, \texttt{p2\_fig1\_new.py}\\
All steps, checks and limits & \texttt{run\_all.sh}, \texttt{compare\_results.py}, \texttt{audit\_numbers.py}, \texttt{README\_REPRODUCE.md}\\\bottomrule
\end{tabularx}
\end{table}
\end{document}
